%=========================== RESUMO ===========================
\begingroup
\begin{center}{\bfseries RESUMO}\end{center}
\vspace{0.3cm}

% ESCREVER POR ÚLTIMO, depois dos resultados. O que está aqui já
% segue a fórmula de cinco movimentos: contexto, impasse, o que foi
% feito, como foi avaliado, o que se achou. Os dois últimos só podem
% ser fechados após a coleta.

\noindent Fobias específicas a animais impõem limitações cotidianas a
parcela expressiva da população, e a superação desse tipo de medo depende
do contato progressivo com o estímulo temido --- que é justamente aquilo
que a pessoa evita. Ambientes virtuais oferecem uma via para contornar
esse impasse, permitindo graduar o estímulo enquanto o controle permanece
com quem é exposto. Este trabalho apresenta o desenvolvimento e a
validação de um protótipo de Realidade Virtual executável em navegador,
sem instalação, que simula cinco animais associados a fobias --- cobra,
rato, barata, aranha e sapo --- cada um no ambiente em que seria
encontrado, com a exposição organizada em cinco níveis progressivos nos
quais variam a distância, a quantidade de exemplares, o movimento e a
possibilidade de interação direta. A aplicação registra automaticamente
medidas de comportamento durante o uso, entre elas a menor distância
alcançada, a proporção do tempo com o estímulo no campo de visão e o
nível máximo atingido, confrontadas com o grau de medo declarado
previamente pelo participante. [COMPLETAR: a validação foi conduzida
com N participantes dos cursos de X e Y.] [COMPLETAR: o que os dados
indicaram.] Ressalta-se que o sistema constitui protótipo acadêmico, não
se destinando a tratamento ou diagnóstico.

\vspace{0.4cm}
\noindent\textbf{Palavras-chave:} Realidade Virtual; Fobias específicas;
Exposição gradual; WebXR; Usabilidade.

\vspace{1.0cm}
\begin{center}{\bfseries ABSTRACT}\end{center}
\vspace{0.3cm}

% Traduzir depois que o português estiver fechado.
\noindent Specific animal phobias impose everyday limitations on a
significant portion of the population, and overcoming such fears depends
on progressive contact with the feared stimulus --- precisely what the
person avoids. Virtual environments offer a way around this impasse,
allowing the stimulus to be graded while control remains with the person
being exposed. This work presents the development and validation of a
browser-based Virtual Reality prototype, requiring no installation, that
simulates five phobia-associated animals --- snake, rat, cockroach,
spider and frog --- each in the environment where it would be
encountered, with exposure organized into five progressive levels varying
distance, number of specimens, movement and the possibility of direct
interaction. The application automatically records behavioural measures
during use, including the closest distance reached, the proportion of
time with the stimulus in the field of view and the highest level
attained, which are compared against the degree of fear previously
declared by the participant. [TO COMPLETE.] The system is an academic
prototype and is not intended for treatment or diagnosis.

\vspace{0.4cm}
\noindent\textbf{Keywords:} Virtual Reality; Specific phobias; Gradual
exposure; WebXR; Usability.
\endgroup
\newpage
